problem:  
Let \((M^n,g)\) be a smooth compact Riemannian manifold with \(\operatorname{Ric}_g \ge (n-1)k\, g\) for some real \(k\), and let \(I_M(v)\) denote its isoperimetric profile function.  
It is known (Bayle, Cavalletti–Mondino, etc.) that \(I_M\) satisfies a second‑order differential inequality of the form  
\[
I_M''(v) + \Phi_{k,n}(I_M'(v), I_M(v)) \le 0,
\]
interpreted in the viscosity sense, where equality holds for the constant‑curvature model \(I_{k,n}\).

**Conjecture (Second‑order equality at one volume implies local rigidity):**  
Suppose that for some nontrivial \(v_0\in(0,\operatorname{Vol}(M,g))\),
\[
I_M(v_0)=I_{k,n}(v_0), \quad I_M'(v_0)=I_{k,n}'(v_0), \quad 
I_M''(v_0)=I_{k,n}''(v_0),
\]
with all derivatives interpreted in the viscosity sense.  
Does this *second‑order contact* at a single volume force local constant‑curvature structure—that is, must the corresponding isoperimetric region \(\Omega_{v_0}\) have a neighborhood in \(M\) isometric to a geodesic ball in the model space form \(M_k^n\)?  

Equivalently: does pointwise equality up to second order in the differential inequality for the isoperimetric profile propagate to local metric rigidity under a lower Ricci bound?

---

Why is it a "good" problem:  

1. **Mathematically well‑posed and precise:**  
   Unlike earlier formulations, all quantities in the conjecture are rigorously defined in the viscosity sense, using a known differential inequality for \(I_M\). This ensures logical and analytic clarity.

2. **Minimalistic and subtle hypothesis:**  
   The assumption demands equality not on an interval but merely *second‑order contact at a single value* of \(v\). Such infinitesimal coincidence is the sharp threshold where analytical curvature comparison might begin to rigidify geometry—making the question both delicate and genuinely unexplored.

3. **Deep analytical geometry interaction:**  
   The conjecture lies at the convergence of PDE and geometry: the isoperimetric profile satisfies a nonlinear, curvature‑dependent PDE, and the question is whether equality at one point in this PDE sense determines the local Riemannian metric, analogous to equality cases in classical Laplacian or Hessian comparison theorems.

4. **Potential for new rigidity mechanisms:**  
   A positive answer would uncover an unprecedented propagation principle for equality of geometric inequalities, while a counterexample would refine the understanding of how lower Ricci curvature bounds control metric and measure structures. Either outcome would be significant.

5. **Extreme simplicity concealing profound depth:**  
   The problem concerns one simple scalar function \(I_M(v)\) and its first two derivatives, yet resolution would touch the foundations of comparison geometry, PDE rigidity, and isoperimetric regularity theory.

Hence, this question is an elegant and challenging research problem—simple to state, rigorously grounded, deeply tied to central themes of modern geometric analysis, and potentially transformative if resolved.